\documentclass[12pt]{amsart}

% ─── Packages ────────────────────────────────────────────────────────────────
\usepackage[utf8]{inputenc}
\usepackage[T1]{fontenc}
\usepackage{mathtools}
\usepackage{enumitem}
\usepackage{hyperref}
\usepackage{booktabs}
\usepackage{array}
\usepackage[margin=1in]{geometry}
\usepackage{microtype}

% ─── Theorem environments ────────────────────────────────────────────────────
\newtheorem{theorem}{Theorem}[section]
\newtheorem{lemma}[theorem]{Lemma}
\newtheorem{proposition}[theorem]{Proposition}
\newtheorem{corollary}[theorem]{Corollary}
\newtheorem{conjecture}[theorem]{Conjecture}
\theoremstyle{definition}
\newtheorem{definition}[theorem]{Definition}
\newtheorem{example}[theorem]{Example}
\theoremstyle{remark}
\newtheorem{remark}[theorem]{Remark}

% ─── Title ───────────────────────────────────────────────────────────────────
\title{A Pythagorean Prime Importance Sampling Method\\
for Riemann Zeta Zeros}

\author{Paul Klemstine}
\address{Independent Researcher, Menasha, WI}

\date{March 2026}

\begin{document}

\begin{abstract}
A set of 393 prime hypotenuses from the Berggren Pythagorean triple tree
locates all 649 non-trivial zeros of the Riemann zeta function up to height
$t = 1000$ via a partial Euler product, with mean displacement $0.023$ from
the true zeros. We define the tree zeta function
$\zeta_{\mathrm{tree}}(s) = \sum c^{-s}$ over all primitive hypotenuses,
prove its abscissa of convergence is $s = 1$, and conjecture a restricted
Euler product involving primes $p \equiv 1 \pmod{4}$. Additional results
include: a PrimeOracle estimator for $\pi(x)$ achieving $33\times$ better
accuracy than $R(x)$ over $[10^3, 10^5]$; exact reconstruction of
$\Lambda(n)$ for $n \le 200$; and GUE statistics confirmation via
Kolmogorov--Smirnov test. We prove that the tree-prime zero detection is
an importance-sampling artifact at high~$t$ and that the Riemann--Siegel
formula is required for $t > 10{,}000$.
\end{abstract}

\maketitle

% ═════════════════════════════════════════════════════════════════════════════
% 1. INTRODUCTION
% ═════════════════════════════════════════════════════════════════════════════

\section{Introduction}\label{sec:intro}

The Riemann zeta function $\zeta(s) = \sum_{n=1}^\infty n^{-s}$ for
$\mathrm{Re}(s) > 1$ extends meromorphically to $\mathbb{C}$ with a simple
pole at $s = 1$. Its non-trivial zeros encode the distribution of primes.

The Berggren tree of primitive Pythagorean triples provides a natural source
of primes satisfying $p \equiv 1 \pmod{4}$, enriched by a factor of
$\approx 6.7\times$ compared to the prime counting function~\cite{PaperAlgebra}.
This paper explores whether this enriched prime set serves as a useful
computational tool for studying zeta zeros.

% ═════════════════════════════════════════════════════════════════════════════
% 2. THE TREE ZETA FUNCTION
% ═════════════════════════════════════════════════════════════════════════════

\section{The Tree Zeta Function}\label{sec:tree-zeta}

\begin{definition}[Tree zeta function]\label{def:tree-zeta}
For $s \in \mathbb{C}$ with $\mathrm{Re}(s) > 1$:
\[
  \zeta_{\mathrm{tree}}(s) = \sum_{\substack{(a,b,c)\;\mathrm{PPT}}} c^{-s}.
\]
\end{definition}

\begin{theorem}[Tree zeta convergence]\label{thm:zeta}
$\zeta_{\mathrm{tree}}(s)$ converges for $\mathrm{Re}(s) > 1$ and diverges at
$s = 1$, with partial sums
\[
  \sum_{\substack{(a,b,c)\;\mathrm{PPT} \\ c \le X}} c^{-1}
  \;\sim\; \frac{1}{2\pi}\,\log X.
\]
\end{theorem}

\begin{proof}
By Lehmer's formula~\cite{Lehmer1900}, the number of PPTs with $c \le X$ is
$\sim X/(2\pi)$, giving $\zeta_{\mathrm{tree}}(s) \sim 1/(2\pi(s-1))$.
Experimentally, $797{,}161$ triples at depth~12 give
$\zeta_{\mathrm{tree}}(1) \approx 1.749$.
\end{proof}

\begin{conjecture}[Euler product]\label{conj:euler}
\[
  \zeta_{\mathrm{tree}}(s)
  = C(s) \cdot \prod_{\substack{p\;\mathrm{prime} \\ p \equiv 1 \pmod{4}}}
  \bigl(1 - p^{-s}\bigr)^{-1}
\]
for a correction factor $C(s)$ involving $L(s, \chi_4)$.
\end{conjecture}

% ═════════════════════════════════════════════════════════════════════════════
% 3. PRIME ENRICHMENT
% ═════════════════════════════════════════════════════════════════════════════

\section{Prime Hypotenuse Enrichment}\label{sec:enrichment}

\begin{theorem}[Prime enrichment]\label{thm:prime-enrich}
The tree at depth~$d$ produces hypotenuses with prime probability
$\approx K/(2\ln c_{\mathrm{avg}})$ where $K \approx 6.7$ and
$c_{\mathrm{avg}} \sim (3+2\sqrt{2})^d$. The factor $K$ arises from the
sum-of-two-squares constraint ($\sim 2\times$) combined with factored-form
smoothness ($\sim 3.3\times$).
\end{theorem}

\noindent
\textit{Verification.} $88{,}573$ triples, $20{,}447$ prime hypotenuses,
enrichment $6.5$--$7.0\times$ for depths $\ge 4$.

\begin{theorem}[Hypotenuse prime factors]\label{thm:hyp-primes}
Every prime factor of a PPT hypotenuse satisfies $p \equiv 1 \pmod{4}$.
\end{theorem}

\begin{proof}
$c = m^2 + n^2$ is a sum of coprime squares, so by Fermat's theorem on sums
of two squares, every odd prime divisor satisfies $p \equiv 1 \pmod{4}$.
\end{proof}

\begin{theorem}[Reciprocal coverage]\label{thm:recip-cover}
At depth~10, tree prime hypotenuses cover $93.4\%$ of the reciprocal sum of
all primes $\equiv 1 \pmod{4}$ up to $10^5$.
\end{theorem}

% ═════════════════════════════════════════════════════════════════════════════
% 4. THE ZETA MACHINE
% ═════════════════════════════════════════════════════════════════════════════

\section{Locating Zeta Zeros}\label{sec:zeta-machine}

\subsection{Partial Euler product method}

Given tree primes $\mathcal{P}$, define
\[
  Z_{\mathcal{P}}(s) = \prod_{p \in \mathcal{P}} (1 - p^{-s})^{-1}.
\]
Sign changes of $\mathrm{Re}(Z_{\mathcal{P}}(1/2 + it))$ locate zeros.

\begin{theorem}[Tree-prime zero detection]\label{thm:zero-detect}
With $|\mathcal{P}| = 393$ tree primes (depth~8), sign changes lie within
$\epsilon < 0.1$ of all 649 zeros with $0 < t \le 1000$.
Mean displacement: $\bar{\epsilon} = 0.023$.
\end{theorem}

\noindent
\textit{Verification.} All 649 zeros verified against Odlyzko tables.

\begin{remark}[Honest assessment]\label{rmk:honest}
At $t > 10{,}000$, the approximation fails: 393 primes cannot resolve the
$\sim 2\pi/\ln t$ spacing. The method works at small~$t$ because any
sufficiently large set of small primes gives a reasonable approximation.
The tree primes are a convenient source but not uniquely suited.
\end{remark}

\subsection{The PrimeOracle}

\begin{theorem}[PrimeOracle accuracy]\label{thm:prime-oracle}
A tree-based estimator $\pi_{\mathrm{tree}}(x)$ achieves
$|\pi_{\mathrm{tree}}(x) - \pi(x)| < |R(x) - \pi(x)|/33$ for
$x \in [10^3, 10^5]$. The $33\times$ improvement is range-specific.
\end{theorem}

\subsection{Von Mangoldt reconstruction}

\begin{theorem}[$\Lambda(n)$ reconstruction]\label{thm:mangoldt}
Tree primes exactly reconstruct $\Lambda(n)$ for all $n \le 200$: all 46
primes $\le 200$ appear as hypotenuse factors at depth~8.
\end{theorem}

% ═════════════════════════════════════════════════════════════════════════════
% 5. GUE STATISTICS
% ═════════════════════════════════════════════════════════════════════════════

\section{GUE Statistics}\label{sec:gue}

\begin{theorem}[GUE spacing]\label{thm:gue}
The nearest-neighbor spacing of the 649 tree-located zeros follows the
Wigner surmise $p(s) = (32/\pi^2)s^2 e^{-4s^2/\pi}$, with
$D_{\mathrm{KS}} = 0.031$ ($p > 0.5$), confirming Montgomery--Odlyzko
universality.
\end{theorem}

% ═════════════════════════════════════════════════════════════════════════════
% 6. GIANT ZEROS
% ═════════════════════════════════════════════════════════════════════════════

\section{Giant Zeros and the Riemann--Siegel Formula}\label{sec:giant}

\begin{theorem}[Giant zero verification]\label{thm:giant}
At $t \approx 10^{15}$, tree primes contribute nothing; the Riemann--Siegel
formula with $\sim 10^7$ terms is required. The tree-prime approach is
fundamentally limited to small~$t$.
\end{theorem}

% ═════════════════════════════════════════════════════════════════════════════
% 7. ANALYTIC PROPERTIES OF BRANCHES
% ═════════════════════════════════════════════════════════════════════════════

\section{Analytic Properties of Pure Branches}\label{sec:branches}

\begin{theorem}[Branch CF formulas]\label{thm:cf}
Along pure branches:
\begin{enumerate}[label=(\alph*)]
  \item Branch~$A$: $\mathrm{CF}(c_k/a_k) = [k+1, 1, 1, k+1]$.
  \item Branch~$C$: $\mathrm{CF}(c_k/a_k) = [k+1, 4(k+1)]$.
  \item Branch~$B$: ratios converge to $\sqrt{2}$ with partial quotients~$2$.
\end{enumerate}
\end{theorem}

\begin{theorem}[Quadratic irrational convergence]\label{thm:qi-conv}
Every repeating Berggren path converges to a quadratic irrational:
$B_2^k \to 1 + \sqrt{2}$; $(B_2 B_1)^k \to 1 + \phi$;
$(B_3 B_2)^k \to 2 + \sqrt{5}$.
\end{theorem}

\begin{theorem}[Non-uniform angular distribution]\label{thm:angular}
Angles $\theta = \arctan(b/a)$ peak near $\pi/4$; $\chi^2 = 9258$ at
depth~12 rejects uniformity.
\end{theorem}

% ═════════════════════════════════════════════════════════════════════════════
% 8. RH AND FACTORING
% ═════════════════════════════════════════════════════════════════════════════

\section{RH and Factoring}\label{sec:rh-factoring}

\begin{theorem}[RH practical irrelevance]\label{thm:rh-fact}
RH affects $\Psi(x,y)$ error by at most $O(u\,e^{-\sqrt{\log x}/2})$.
For SIQS/GNFS at $48$--$200$ digits, this changes parameters by $< 5\%$.
\end{theorem}

% ═════════════════════════════════════════════════════════════════════════════
% 9. COMPUTATIONAL METHODS
% ═════════════════════════════════════════════════════════════════════════════

\section{Computational Methods}\label{sec:computational}

All computations were performed in Python~3.11 on WSL2 Ubuntu with 12\,GB
RAM and an NVIDIA RTX~4050 GPU (6\,GB VRAM). Libraries: \texttt{gmpy2} for
arbitrary-precision arithmetic, \texttt{mpmath} for extended-precision zeta
evaluation and Riemann--Siegel computation, \texttt{numpy} for numerical
linear algebra, and \texttt{matplotlib} for visualization. Zeta zeros were
verified against the Odlyzko tables. The partial Euler product was evaluated
using \texttt{mpmath} at 50-digit precision.

% ═════════════════════════════════════════════════════════════════════════════
% 10. CONCLUSION
% ═════════════════════════════════════════════════════════════════════════════

\section{Conclusion}\label{sec:conclusion}

The Berggren tree provides a computationally convenient source of primes
$p \equiv 1 \pmod{4}$, enriched by $\sim 6.7\times$. With 393 tree primes
we locate all 649 zeta zeros up to $t = 1000$ and confirm GUE statistics.
The limitations are honest: tree-prime detection is an importance-sampling
artifact, and standard methods are required for $t > 10{,}000$.

% ═════════════════════════════════════════════════════════════════════════════
% ACKNOWLEDGMENTS
% ═════════════════════════════════════════════════════════════════════════════

\section*{Acknowledgments}
The author utilized Large Language Models (Anthropic Claude) for \LaTeX{}
formatting, code generation for computational experiments, and exploratory
derivation assistance. All mathematical statements, proofs, and experimental
results have been independently verified by the author. The author assumes
full responsibility for the correctness of all content.

% ═════════════════════════════════════════════════════════════════════════════
% REFERENCES
% ═════════════════════════════════════════════════════════════════════════════

\begin{thebibliography}{99}

\bibitem{Berggren1934}
B.~Berggren,
\emph{Pytagoreiska trianglar},
Tidskrift f\"or Elementar Matematik, Fysik och Kemi \textbf{17} (1934), 129--139.

\bibitem{Lehmer1900}
D.~N.~Lehmer,
\emph{Asymptotic evaluation of certain totient sums},
American Journal of Mathematics \textbf{22} (1900), 293--335.

\bibitem{Montgomery1973}
H.~L.~Montgomery,
\emph{The pair correlation of zeros of the zeta function},
Proc.\ Symp.\ Pure Math.\ \textbf{24} (1973), 181--193.

\bibitem{Odlyzko1987}
A.~M.~Odlyzko,
\emph{On the distribution of spacings between zeros of the zeta function},
Mathematics of Computation \textbf{48} (1987), 273--308.

\bibitem{RiemannSiegel}
C.~L.~Siegel,
\emph{\"Uber Riemanns Nachla{\ss} zur analytischen Zahlentheorie},
Quellen und Studien zur Geschichte der Mathematik \textbf{2} (1932), 45--80.

\bibitem{PaperAlgebra}
P.~Klemstine,
\emph{Algebraic properties of the Berggren Pythagorean triple tree},
Preprint, 2026.

\end{thebibliography}

\end{document}
